\documentclass[../../main.tex]{subfiles}

\graphicspath{{../../image/}{../../figure/}}

\begin{document}

% 原书物理页 9，印刷页码 X--XI。

\chapter*{致谢}
\phantomsection\label{chapter:front-acknowledgments}
\pdfbookmark[1]{致谢}{bm-front-acknowledgments}
\markboth{致谢}{致谢}

我的引路人 Jerry Ericksen 与 Clifford Truesdell 把我领入连续介质物理——它既是一门活生生的科学主题，又是一种具有迷人历史的形式化数学结构。我相信，这两种视角都以某种方式反映在本书之中。

我感谢许多科学家——老师、同事与学生一视同仁——他们在过去五十年间帮助我学习连续介质物理与双曲守恒律理论。由于无法在此一一列举他们的姓名，请允许我特别指出 Stu Antman、John Ball、Alberto Bressan、Gui-Qiang Chen、Cleopatra Christoforou、Bernie Coleman、Ron DiPerna、Hermano Frid、Jim Glimm、Jim Greenberg、Mort Gurtin、Ling Hsiao、Barbara Keyfitz、Peter Lax、Philippe LeFloch、Tai-Ping Liu、Andy Majda、Piero Marcati、Ingo Müller、Walter Noll、Jim Serrin、Denis Serre、Marshall Slemrod、Joel Smoller、Luc Tartar、Konstantina Trivisa、Thanos Tzavaras、Dehua Wang 与 Zhouping Xin，他们同样以友谊相待。特别地，Denis Serre 的不断鼓励帮助我把这项艰巨的工作推向完成。

我深深感谢 Jennifer Curtiss Gage，她就写作风格问题所给出的坚定而睿智的建议实在不可或缺。

我尤其感念 Mengrui Ni，他为录入一部变化不定的手稿付出了西西弗斯式的辛劳。我也感谢 Changqing (Peter) Hu、Vasileios Symeonidis 与 Mihalis Dafermos 根据我的粗疏草图绘制插图。我同样感谢 Brown University 的 Stephanie Han，以及 Springer 数学编辑部（Springer Mathematics Editorial Department）的 Catriona Byrne 与 Marina Reizakis 的友好合作。

\end{document}